%
% This file is part of the project of
% National Cheng Kung University (NCKU) Thesis/Dissertation Template in LaTex.
% This project is hold at
%     <https://github.com/wengan-li/ncku-thesis-template-latex>
% by Wen-Gan Li.
%
% This project is distributed in the hope of usefuling to someone,
% you can redistribute it and/or modify it under the terms of the
% Attribution-NonCommercial-ShareAlike 4.0 International.
%
% You should have received a copy of the
% Attribution-NonCommercial-ShareAlike 4.0 International
% along with this project.
% If not, see <http://creativecommons.org/licenses/by-nc-sa/4.0/legalcode.txt>.
%
% Please feel free to fork it, modify it, and try it.
% Have fun !!!
%

% ----------------------------------------------------------------------------

\DeclareDocumentCommand{\SetTableCaptionAndLabel}{+m +m}
{%
  \SetFloatCaptionAndLabel{#1}{#2}%
} % End of \DeclareDocumentCommand{}

\DeclareDocumentCommand{\SetTableCaptionStarAndLabel}{+m +m}
{%
  \ifthenelse{\equal{#1}{\empty}}{}%
  {%
    \caption*{#1}%
    \NCKUPrivateFreezeCurrentLabelName{#1}%
    \ifthenelse{\equal{#2}{\empty}}{}{\label{#2}}%
  }%
} % End of \DeclareDocumentCommand{}

\newcommand{\DisplayTableContent}[5]
{%
  \NCKUPrivateDisplayFramedFloatContent{#4}{%
    \ifthenelse{\equal{#1}{0.0}}%
    {%
      \makebox[\textwidth]%
      {%
        \setlength{\tabcolsep}{#2}%
        \renewcommand{\arraystretch}{#3}%
        #5%
      }%
    }{%
      \makebox[\textwidth]%
      {%
        \resizebox{#1\paperwidth}{!}%
        {%
          \setlength{\tabcolsep}{#2}%
          \renewcommand{\arraystretch}{#3}%
          #5%
        }%
      }%
    }%
  }%
} % End of \newcommand{}

\ExplSyntaxOn
\keys_define:nn { ncku / insert-table }
  {
    scale        .tl_set_e:N = \TmpValueScale,
    nomtitle     .tl_set_e:N = \TmpValueNomTitle,
    caption      .tl_set_e:N = \TmpValueCaption,
    label        .tl_set_e:N = \TmpValueLabel,
    pos          .tl_set_e:N = \TmpValuePosition,
    tabcolsep    .tl_set_e:N = \TmpValueTabColSep,
    arraystretch .tl_set_e:N = \TmpValueArrayStretch,
    opacity      .tl_set_e:N = \TmpValueOpacity,
  }

% Reset every value on each call, then apply the caller's overrides. The public
% optional argument uses the literal token \empty when omitted.
\cs_new_protected:Npn \ncku_insert_table_keys_set:n #1
  {
    \keys_set:nn { ncku / insert-table }
      {
        scale = 0.0,
        nomtitle = \empty,
        caption = \empty,
        label = \empty,
        pos = top,
        tabcolsep = 6pt,
        arraystretch = 1,
        opacity = 0.4,
      }
    \tl_if_eq:nnF {#1} {\empty}
      { \keys_set:nn { ncku / insert-table } {#1} }
  }
\cs_new_eq:NN \NCKUPrivateSetInsertTableKeys
  \ncku_insert_table_keys_set:n
\ExplSyntaxOff

\newcommand{\InsertTable}[2][\empty]
{%
  % Parse the input
  \NCKUPrivateSetInsertTableKeys{#1}%
  %
  \begin{table}[H]%
  %
    \ifthenelse{\equal{\TmpValuePosition}{top}}%
    {%
      \ifthenelse{\equal{\TmpValueNomTitle}{\empty}}%
        {\SetTableCaptionAndLabel{\TmpValueCaption}{\TmpValueLabel}}%
        {\SetTableCaptionStarAndLabel{\TmpValueNomTitle}{\TmpValueLabel}}%
      \DisplayTableContent{%
        \TmpValueScale}{\TmpValueTabColSep}{%
        \TmpValueArrayStretch}{\TmpValueOpacity}{#2}%
    }%
    {%
      \DisplayTableContent{%
        \TmpValueScale}{\TmpValueTabColSep}{%
        \TmpValueArrayStretch}{\TmpValueOpacity}{#2}%
      \ifthenelse{\equal{\TmpValueNomTitle}{\empty}}%
        {\SetTableCaptionAndLabel{\TmpValueCaption}{\TmpValueLabel}}%
        {\SetTableCaptionStarAndLabel{\TmpValueNomTitle}{\TmpValueLabel}}%
    }%
  \end{table}%
} % End of \newcommand{}

% -----------------------------------------------------------------

\newcolumntype{L}[1]{>{\raggedright\let\newline\\\arraybackslash\hspace{0pt}}m{#1}}
\newcolumntype{C}[1]{>{\centering\let\newline\\\arraybackslash\hspace{0pt}}m{#1}}
\newcolumntype{R}[1]{>{\raggedleft\let\newline\\\arraybackslash\hspace{0pt}}m{#1}}

% -----------------------------------------------------------------

\def\ValueTableNameDefault{Table}
\def\ValueTableNameCustom{Table} % Default
\def\UseTableNameDefault{%
  \renewcommand{\tablename}{\ValueTableNameDefault}}
\def\UseTableNameCustom{%
  \renewcommand{\tablename}{\ValueTableNameCustom}}
\newcommand{\SetCustomTableName}[1]{%
  \renewcommand{\ValueTableNameCustom}{#1}}

\UseTableNameCustom % Default

% Default style
\newcommand{\UseTableCaptionDefaultStyle}
{%
  \captionsetup*[table]{labelfont=bf, textfont=normalfont}%
} % End of \newcommand{}

% Style for Extended Abstract
\newcommand{\UseTableCaptionExtendedAbstractStyle}
{%
  \captionsetup*[table]{font=bf}%
  \renewcommand{\thetable}{\arabic{table}}%
} % End of \newcommand{}

\UseTableCaptionDefaultStyle % Default

% -----------------------------------------------------------------
